\documentclass[11pt]{article}
\usepackage{amsmath,amssymb}
\usepackage[hidelinks]{hyperref}
\title{Inferring a handoff from timestamps and incomplete records}
\date{}
\begin{document}
\maketitle

With two source records and two destination records, every pairing
can look possible when checked separately. A shared clock offset
can nevertheless force the true pairing. In the model below that
is enough to retire an old duty, but only if both transfers are
known to be represented. Allowing even one missing transfer
can make it unsafe to release the same reserve.

\section{Source law}
Two read streams use their own clock origins. Event \(i\) is
logged at source time \(s_i\) and event \(j\) at destination
reference time \(r_j\). All source records in the episode
use one unknown common displacement \(x\), with independently
validated limits. A physically admissible candidate transfer
\(ij\) has elapsed time in \([\ell_{ij},u_{ij}]\). Transfer
eligibility must come from the actual transport process or be
declared as an abstraction. For this result there is no
additional shared travel-capacity constraint or phase-dependent
service law beyond uniqueness and the stated pair intervals.
The exact pair condition is
\[
        x \in [\,r_j-s_i-u_{ij},\ r_j-s_i-\ell_{ij}\,].
\]
Intersect every such interval with the admitted calibration.
Using a different offset for each pair is not the same source.

A complete explanation uses one-to-one assignments of source
records to destination records. For an incomplete observation
interface, admit a matching with cardinality \(k\) in a proven
window \([h,K]\). For instance, a separate trustworthy
event counter may establish that at least \(h\) destination
records came from the listed distinct sources. This is not
the number of lines in the log. It does not follow from
absence of an error or enough unrelated pulses. Duplicated
or reordered messages, missing event types, and physical
non-occurrence must be resolved or admitted separately.
If no evidence establishes a positive lower count, it is
not safe to promote the records to a perfect matching.

Suppose the reserve exposure in a matching \(M\) is \(w_{ij}\)
for a matched pair \(ij\), \(a_i\) for an unmatched source,
and \(b_j\) for an unmatched destination. The score is
\[
 \sum_i a_i+\sum_j b_j+
     \sum_{ij\in M}(w_{ij}-a_i-b_j).
\]
These scores refer to the declared old obligations, not to
timestamps alone. For example an unresolved source may still
owe a physical duty even if the only connection to its destination
is missing. A confirmed output at the wrong destination need
not settle this particular duty either. If source risks vary
with offset or some incident history other than assignment,
they must be included in the model; the simple constant-weight
formula does not establish them.

\section{Exact extrema and proof witnesses}
For a fixed common offset the problem is weighted bipartite
matching with a cardinality window. For exactly \(k\) matches
between \(n\) source and \(m\) destination records, add \(m-k\)
dummy sources attached to all real destinations and \(n-k\)
dummy destinations attached to all real sources. There are no
dummy--dummy edges. Any perfect assignment of the enlarged
graph must match \(m-k\) real destinations to dummy sources,
so exactly \(k\) real pairs remain. Conversely every size-\(k\)
matching extends to a perfect assignment by connecting the
unmatched records to their corresponding dummy vertices.
Use the adjusted weight \(w_{ij}-a_i-b_j\) on a real edge
and zero on a dummy edge. Trying the declared sizes and the
two cost signs gives the exact minimum and maximum of the
exposure. No assertion is made that every intermediate score
is attained.

For one common scalar offset, the eligible graph is constant
on every open cell of the pair-interval endpoints. At an
endpoint the graph can differ from both adjacent open cells.
Enumerate the endpoints and one rational interior point per cell.
If the calibration extends outside the pair endpoints, no
real edge exists there and only an allowed empty matching matters.
For a cell and cardinality that has no assignment, a Hall
deficiency is a certificate. For one with an assignment, the
checker uses feasible integer row and column duals with
equality at an actual matching; summing dual inequalities
proves optimality without trusting the floating cost solver.
It checks both cost signs and the full set of offset cells
and sizes for the fixed fixture. This certifies the output for
the supplied interval and exposure model, not how to derive
that model from measurements. Several coupled clock dimensions,
unbounded drift or interaction between simultaneous transfers
can require a different feasibility solver.

\section{A release decision that changes under missing records}
Take \(s=(0,2)\), \(r=(0,2)\), \(x\in[-3,3]\) and allowable
elapsed time \([0,1]\) on every pair. The two straight pairs
both agree on \(x\in[-1,0]\). The cross pairs individually
allow \(x\in[1,2]\) and \(x\in[-3,-2]\), respectively.
They cannot both hold in one history. With two actual
matched pairs the straight assignment is necessary. If
source zero carries one named old duty discharged only by
arrival at destination zero, the exact worst residual
exposure is zero. A graph made just from marginally possible
edges would also admit the cross assignment and could keep
the reserve locked unnecessarily.

When just one transfer is guaranteed, the straight
assignment remains possible but a source or destination can
be unmatched, and a cross pair can occur alone. The same
reserve now has worst exposure one. Authorizing its reuse
from the perfect-assignment result would be unsafe.
The distinction applies even with exactly the same two
sets of timestamps and the same clock calibration.
When a deciding actor has yet to receive evidence of the
count, documented archival receipt does not let that actor
use the tighter source. An online controller needs the
right to read those events by its actual commitment,
and a genuine certification of transfer, not the freedom
to select the most favorable compatible history.

\section{Evidence status}
\href{https://ftp.eecs.umich.edu/~pettie/matching/Kuhn-hungarian-assignment.pdf}{Kuhn's
assignment duality} is standard. The clock-cell method and
dummy reduction are applications of existing machinery;
same-input enumeration ties. Five Lean statements check the
time algebra, marginal-versus-joint fixture, and dummy
cardinality step. The broader sweep and weights are checked
by separate finite enumeration and checked certificates.
The externally published \href{https://zenodo.org/records/4034867}{CONTEXT
release} motivates the question, since its streams use
approximately aligned clocks and can have missing or
out-of-order observations. No numerical data from that
release are used and no physical timing or completeness
bound is inferred from it. This is an I1 calibration
case, not an independent transfer or data-derived estimate.
\end{document}
